\begin{frame}
\titlepage
\note{Tuesday is 75 minutes. Keep the rhythm: predict, run or fill the evidence, explain what survived, and name the boundary.}
\end{frame}

\begin{frame}{The question for today}
\princ{MODEL SCOPE}
\begin{block}{Week 14 question}
Is this transformation equivalent under every assignment, and what can this state model not establish?
\end{block}
\begin{itemize}
\item A Boolean expression says what is true for one assignment.
\item A finite-state machine says how a bounded model changes as inputs arrive.
\item Both are useful only when their assumptions are visible.
\end{itemize}
\predict{What evidence would make “these expressions are equivalent” a justified claim?}
\see{every relevant assignment is checked, not just a favorable example.}
\note{Connect the two halves: exhaustive rows for Boolean claims, explicit traces and transitions for state claims.}
\end{frame}

\begin{frame}{Our checking rhythm}
\princ{PREDICT--RUN--EXPLAIN}
\begin{enumerate}
\item Predict the row, witness, trace, or missing case.
\item Run a small, visible check or fill it by hand.
\item Explain exactly what the evidence establishes.
\item State one boundary: what was not represented or checked?
\end{enumerate}
\begin{block}{A useful rule}
A fluent answer is a proposal. A complete check is evidence.
\end{block}
\see{each section ends with a row, counterexample, trace, or gap.}
\end{frame}

\begin{frame}{Words we will use}
\begin{tabular}{@{}p{3.0cm}p{9.8cm}@{}}
\textbf{Boolean variable} & A symbol that receives either true or false.\\[2pt]
\textbf{assignment} & The chosen true/false value for every variable.\\[2pt]
\textbf{operator} & A rule such as AND, OR, or NOT that combines or changes values.\\[2pt]
\textbf{expression} & A written combination of variables and operators.\\[2pt]
\textbf{truth table} & A table with one row for every assignment.\\[2pt]
\textbf{equivalent} & Producing the same result for every assignment in scope.\\[2pt]
\textbf{counterexample} & One in-scope case that disproves a universal claim.\\[2pt]
\textbf{circuit} & A picture of gates connected to represent an expression.\\[2pt]
\textbf{gate} & A circuit component that performs one Boolean operation.\\[2pt]
\textbf{state} & A named condition in a changing system.\\[2pt]
\textbf{transition} & A rule that changes one state to another for an input.\\[2pt]
\end{tabular}
\note{Do not read every line aloud. Point out that the terms are defined before we rely on them.}
\end{frame}

\begin{frame}{The tempting simplification}
\princ{CHECK THE CLAIM}
\begin{columns}[T]
\begin{column}{0.48\textwidth}
\begin{block}{Supplied expression}
\[
F=(P\ \text{OR}\ Q)\ \text{AND}\ (P\ \text{OR}\ \text{NOT }Q)
\]
\end{block}
\end{column}
\begin{column}{0.48\textwidth}
\begin{alertblock}{Confident proposed simplification}
\[
B=P\ \text{OR}\ Q
\]
\end{alertblock}
\end{column}
\end{columns}
\predict{If an AI says $F$ and $B$ are equivalent, what is the smallest check that could break the claim?}
\see{a complete four-row truth table can find a single counterexample.}
\mistake{“They match on the first row” is not the same as “they match on every row.”}
\end{frame}
